Con el ánimo de aclarar de qué va este proyecto, así es como lo veo yo.

1) La novedad esencial y lo que ha cambiado con respecto al año pasado es darnos cuenta de que existe un nicho para la tecnología CRYSP. Aplicación a terapia de protones.

2) La tecnología CRYSP, conjunta varias propuestas: 1) PET monolítico; 2) electrónica custom; 3) Materiales baratos. Ninguna de estas cosas es novedad, la gente lleva años intentando que funcionen los monolíticos (hasta ahora sin mucho éxito) y buscando alternativas al LYSO (hasta ahora solo el BGO para según que cosas). Lo novedoso de CRYSP es la posibilidad de que, al conjuntar materiales baratos y disponibles con electrónica Custom y buena respuesta física tengas un PET barato de buenas prestaciones.

3) Aún así, competir en el mercado libre con los gigantes de PET parece muy difícil. El CsI frío es razonablemente competitivo (como demostramos en el paper de CRYSP) para un caso de PET de 1 m, aunque tiene algunas limitaciones de sensitividad y de timing. El BGO frío no funciona, porque el timing es definitivamente demasiado largo. 

4) El nicho de la terapia de protones cambia las cosas. Aquí necesitas un PET de buena resolución espacial y de energía, gracias al DOI los monolíticos son ideales, el timing no es un problema y la sensitividad es menos problema que en el caso normal. Poner un PET in-room, por otra parte, pide que, el manos el demostrador sea barato y sencillo. De esta manera, nadie de los grandes va a intentar comerte la tostada (tu PET es mucho más barato que el suyo) y esto te da la opción de afianzar tu nicho.

5) Lo más importante para afianzar ese nicho es cerrar un ciclo completo. Construir un demostrador, instalarlo en el centro de terapia y demostrar que les sirve a los médicos. En mi opinión ese es nuestro objetivo ahora. 

6) El demostrador más fácil es de CsI. El mejor es de BGO a 195 K. No hace falta que decidamos hoy, pero sí que tenamos claros los plazos y que no corramos riesgos. El CsI es algo que sabemos hacer y debería ser nuestro defecto hasta que no se demuestre lo contrario. A la vez podemos desarrollar el criostato y si todo va bien, solo hay que reemplazar un cristal por otro. 

7) Lo crucial sin embargo es construir el aparato, interesar a los médicos, ponerlo en marcha, desarrollar la electrónica, afianzar la tecnología. Todo esto nos puede llevar entre dos y 4 años. Tenemos financiación para ello y hay que desarrollar el equipo que lo haga.

8) Como entra la empresa en todo esto hay que hablarlo. Creo que como spinoff que es debe estar ahí desde el principio. La financiación incial seguramente vendrá de ciencia, vía proyecto Faro, EIC, o una combinación de todo. Eso nos dará recorrido sin necesidad de salir como locos a buscar inversores. 

9) Claramente, intentarlo ahora sería prematuro, ya que no tenemos nada. En 2 años podemos tener un PET, el software, el know-how, la electrónica y el nicho. Y hasta entonces tenemos funding. 

10) Yo creo que ese es el camino a seguir. Las soluciones técnicas finales están por ver. Si somos capaces de fabricar un PET y que los médicos lo usen con éxito, creo que tendremos un caso fuerte y mucho prestigio. Como dices tú, hasta ahí puedo leer. 

